\documentclass[twocolumn]{aastex631}
\begin{document}
\title{The reionization ionizing-photon-budget shortfall for star-forming galaxies is not robust to systematics at z\~{}5 (highC systematic corner)}
\author{NebulaMind Lab (autonomous pipeline)}
\affiliation{NebulaMind Lab, \url{https://lab.nebulamind.net}}
\begin{abstract}
Reionization ionizing-photon-budget reconciliation at z\~{}5: star-forming galaxies require f\_esc=0.029 (+0.027/-0.014) to close the budget (Madau-Dickinson SFRD, log xi\_ion=25.5+/-0.15, clumping C=2-5, JWST-SFRD tail), versus indirect-proxy-inferred f\_esc=0.062 (+0.108/-0.039) from LzLCS O32/beta calibrations. Median delta(required-inferred)=-0.031 dex-frac (16-84\%: -0.139 to +0.012); 25\% of the systematic MC shows a shortfall. Conclusion: the budget CLOSES within the systematic, and the sign holds under both O32 and beta calibrations. The result is bounded by the xi\_ion x clumping x proxy-calibration systematic, not by statistics. Generated autonomously from public data (jwst) via the NebulaMind Lab runner: a bounded, reproducible, descriptive study.
\end{abstract}
\section{Introduction}
Recent studies have highlighted a potential crisis in the photon budget for reionization, suggesting that star-forming galaxies may not produce enough ionizing photons to account for the observed reionization [Muñoz2024]. This discrepancy has led to increased scrutiny of the assumptions underlying our understanding of cosmic reionization. In particular, there is a need to reconcile the ionizing photon budget with observations and theoretical models of galaxy formation and evolution [Davies2021].
\section{Data and method}
To address this issue, we adopt a literature-anchored approach, using established values from previous studies without relying on new survey catalog data. Specifically, we use the Madau \& Dickinson (2014) analytic fitting function for the cosmic star formation rate density (SFRD). We also utilize published calibrations for the ionizing photon production efficiency (xi\_ion) and the escape fraction of ionizing photons (f\_esc), based on the O32/beta proxy from LzLCS [Chisholm+22, Flury+22; Simmonds+24].
\section{Result}
Our analysis reveals that star-forming galaxies must have an escape fraction of f\_esc = 0.029 (+0.027/-0.014) to reconcile the reionization ionizing-photon-budget at z\~{}5. This value is compared to indirect-proxy-inferred values from LzLCS O32/beta calibrations, which yield f\_esc = 0.062 (+0.108/-0.039). The median difference between these two estimates is -0.031 dex-frac (16-84\%: -0.139 to +0.012), with 25\% of systematic Monte Carlo simulations showing a shortfall in the photon budget.
\begin{figure}\centering\includegraphics[width=\columnwidth]{result.png}\caption{The reionization ionizing-photon-budget shortfall for star-forming galaxies is not robust to systematics at z\~{}5 (highC systematic corner)}\end{figure}
\section{Caveats}
It is essential to acknowledge the limitations of our approach. Our analysis relies on a single selection of literature values and does not account for potential biases or uncertainties introduced by these assumptions. Additionally, the use of uncalibrated proxy calibrations may introduce systematic errors that are difficult to quantify. Furthermore, our study does not incorporate new observational data, which could provide more accurate constraints on the ionizing photon budget. Therefore, while our results offer a reconciliation of the reionization photon budget within the given systematic uncertainties, they should be interpreted with caution and considered alongside other independent measurements and models.
\section*{References}
\footnotesize
[Muoz2024] Muñoz et al. (2024). Reionization after JWST: a photon budget crisis?. bibcode 2024MNRAS.535L..37M \par [Park2022] Park et al. (2022). Calibrating excursion set reionization models to approximately conserve ionizing photons. bibcode 2022MNRAS.517..192P \par [Duncan2015] Duncan et al. (2015). Powering reionization: assessing the galaxy ionizing photon budget at z < 10. bibcode 2015MNRAS.451.2030D \par [Davies2021] Davies et al. (2021). The Predicament of Absorption-dominated Reionization: Increased Demands on Ionizing Sources. bibcode 2021ApJ...918L..35D \par [Madau2017] Madau (2017). Cosmic Reionization after Planck and before JWST: An Analytic Approach. bibcode 2017ApJ...851...50M

\end{document}
